\documentclass[11pt]{article}

\usepackage[margin=1in]{geometry}
\usepackage{amsmath,amssymb,amsthm}
\usepackage{algorithm}
\usepackage{algpseudocode}
\usepackage{booktabs}
\usepackage[expansion=false]{microtype}
\usepackage{xcolor}
\usepackage[colorlinks=true,linkcolor=blue,citecolor=blue]{hyperref}

\algrenewcommand\algorithmicrequire{\textbf{Input:}}
\algrenewcommand\algorithmicensure{\textbf{Output:}}

\newtheorem{remark}{Remark}
\newtheorem{proposition}{Proposition}

\newcommand{\R}{\mathbb{R}}
\newcommand{\T}{\mathsf{T}}
\newcommand{\Fnorm}[1]{\left\|#1\right\|_F}
\newcommand{\rank}{\operatorname{rank}}

\title{A Blocked Randomized Range-Finder Baseline\\
for Tile Low-Rank GEMM:\\
GPU Implementation Blueprint}
\author{}
\date{}

\begin{document}
\maketitle

\begin{abstract}
We give an implementation blueprint for the update of the output tiles of a tile
low-rank (TLR) matrix product $C\leftarrow\alpha AB+\beta C$ using a blocked
adaptive randomized range finder (ARA), in contrast to the deterministic
sequential-recompression baseline used by libraries such as KBLAS. The
contribution of the reduction index is never materialized as a dense tile, nor
as an explicit running low-rank accumulator recompressed at every contraction
step; it is kept as a \emph{factor list} and applied implicitly through cheap
batched GEMMs. We cost the deterministic baseline first (\S\ref{sec:strategies})
to fix the metric the randomized design is judged against, show how
distributivity turns the whole reduction over the contraction index into a
sketching problem solved with GEMMs rather than $q_k$ separate QR/SVD
recompressions (\S\ref{sec:factorlist}), show how the existing $\beta C_{ij}$
term is folded into the same range finder as one further factor pair, and
derive a workspace-driven rolling schedule that traverses the output one block
row (or block column) at a time. The design targets stock batched GPU
primitives (batched GEMM, SYRK, POTRF, TRSM, and a batched symmetric
eigensolver) with no dependence on routines that form Householder $Q$ factors.
\end{abstract}

\tableofcontents

\section{Setting and notation}
\label{sec:setting}

Let $A\in\R^{M\times K}$ and $B\in\R^{K\times N}$ be stored in TLR form with
tiles of size $b_m\times b_k$ and $b_k\times b_n$, and set $q_m=M/b_m$,
$q_k=K/b_k$, $q_n=N/b_n$. Off-diagonal tiles are stored as outer products,
diagonal tiles densely. For the low-rank tiles,
\[
A_{i\ell}=U^A_{i\ell}\bigl(V^A_{i\ell}\bigr)^{\T},
\qquad
B_{\ell j}=U^B_{\ell j}\bigl(V^B_{\ell j}\bigr)^{\T},
\qquad
C_{ij}=U^C_{ij}\bigl(V^C_{ij}\bigr)^{\T},
\]
The implementation uses compact ASCII spellings of these axis-indexed quantities:
$(b_m,b_k,b_n)\leftrightarrow(\texttt{bm},\texttt{bk},\texttt{bn})$ and
$(q_m,q_k,q_n)\leftrightarrow(\texttt{qm},\texttt{qk},\texttt{qn})$.
Likewise, $(r_A,r_B,r_C)$ are written as `rA`, `rB`, and `rC`; the four
logical factor handles correspond to $U^A,V^A,U^B,V^B$ in that order.
with ranks $r_{A,i\ell}$, $r_{B,\ell j}$, $r_{C,ij}$. Any product of two such
tiles collapses through a small bilinear pairing, the \emph{coupling matrix}
\begin{equation}
\label{eq:coupling}
S_{i\ell j}:=\bigl(V^A_{i\ell}\bigr)^{\T}U^B_{\ell j}
\in\R^{r_{A,i\ell}\times r_{B,\ell j}},
\qquad
A_{i\ell}B_{\ell j}=U^A_{i\ell}S_{i\ell j}(V^B_{\ell j})^{\T},
\end{equation}
independent of how the sum over $\ell$ is subsequently handled --- the two
alternatives are the subject of \S\ref{sec:strategies}. We update one output
tile
\begin{equation}
\label{eq:update}
X_{ij}=\alpha\sum_{\ell=1}^{q_k}A_{i\ell}B_{\ell j}+\beta C_{ij}
\end{equation}
and return it in low-rank form to the Frobenius tolerance
\begin{equation}
\label{eq:tol}
\Fnorm{X_{ij}-U_{ij}V_{ij}^{\T}}\le\tau_{ij},
\qquad
\tau_{ij}=\max\bigl(\varepsilon_{\mathrm{abs}},\,
\varepsilon_{\mathrm{rel}}\Fnorm{X_{ij}}\bigr).
\end{equation}
The absolute floor keeps the criterion well posed when cancellation makes the
updated tile numerically zero. We write
$\rho_i=\sum_\ell r_{A,i\ell}$ for the total left rank of block row $i$,
$\gamma_j=\sum_\ell r_{B,\ell j}$ for the total right rank of block column $j$,
and
$\hat\rho_{ij}=r_{C,ij}[\beta\ne0]+\sum_\ell\min(r_{A,i\ell},r_{B,\ell j})$
for the width of the factor list of \S\ref{sec:factorlist}. All three are known
from rank metadata alone, at no arithmetic cost, and drive the schedule of
\S\ref{sec:schedule}. The chief target of this paper's design is
$\hat\rho_{ij}$, not $b_mb_n$ and not $q_k$: \S\ref{sec:strategies} costs a
deterministic scheme that pays for $q_k$ and a randomized one that does not.

\section{Two strategies for the TLR product}
\label{sec:strategies}

There is more than one correct way to realize \eqref{eq:update} at the tile
level. This section fixes a strong deterministic baseline, in the style of
KBLAS-family TLR GEMM implementations, states its cost precisely, and uses
that cost to motivate the randomized design developed in the rest of the
paper.

\subsection{Deterministic sequential recompression}
\label{sec:kblas}

KBLAS assumes a sum-of-outer-products formulation of \eqref{eq:update}
evaluated by literal, deterministic recompression: the reduction index $\ell$
is processed one step at a time, and after each step the running accumulator
is folded back down to a single low-rank pair before moving on to $\ell+1$.
Seeding $(U,V)\gets(U^C_{ij},\beta V^C_{ij})$ when $\beta\ne0$ (empty
otherwise):

\begin{algorithm}[H]
\caption{Deterministic sequential recompression, one output tile
(KBLAS-style)}
\label{alg:kblas}
\begin{algorithmic}[1]
\Procedure{DeterministicUpdate}{$\{A_{i\ell}\}_{\ell=1}^{q_k},\,
\{B_{\ell j}\}_{\ell=1}^{q_k},\,\beta C_{ij}$}
  \State $(U,V)\gets(U^C_{ij},\,\beta V^C_{ij})$ if $\beta\ne0$, else empty
  \For{$\ell=1$ to $q_k$}
     \Comment{sequential: step $\ell$ needs step $\ell-1$'s recompressed output}
     \State $S_{i\ell j}\gets(V^A_{i\ell})^{\T}U^B_{\ell j}$
        \Comment{\eqref{eq:coupling}, batched over all $(i,j)$}
     \State $U\gets[\,U\ \ U^A_{i\ell}\,]$;\quad
            $V\gets[\,V\ \ \alpha V^B_{\ell j}S_{i\ell j}^{\T}\,]$
        \Comment{concatenate: rank grows by $r_{A,i\ell}$}
     \State $(Q_U,R_U)\gets\textsc{QR}(U)$;\quad
            $(Q_V,R_V)\gets\textsc{QR}(V)$
        \Comment{2 batched QRs, explicit Householder $Q$}
     \State $(W,\Sigma,P)\gets\textsc{SVD}(R_UR_V^{\T})$; truncate to the
            working rank $\hat r$
     \State $U\gets Q_UW(:,1{:}\hat r)$;\quad
            $V\gets Q_VP(:,1{:}\hat r)\Sigma(1{:}\hat r)$
        \Comment{recompress before the next step}
  \EndFor
  \State \Return $(U,V)$
\EndProcedure
\end{algorithmic}
\end{algorithm}

Every line inside the $\ell$-loop is one batched call over all $q_mq_n$ output
tiles; the loop itself is sequential because step $\ell$ needs step $\ell-1$'s
recompressed factors. This costs, per tile, \emph{exactly $q_k$ compressions}
--- batched over the whole grid, that is $q_k$ batched QR-plus-SVD
compressions with batch size $q_mq_n$, i.e.\ $O(q_kq_mq_n)$ small dense
factorizations in total, against $O(q_mq_n)$ for a scheme that compresses
each tile once. Two costs compound this:

\begin{itemize}
\item \textbf{The QR forms an explicit $Q$.} Recompressing
  $[\,U\ \ U^A_{i\ell}\,]$ and $[\,V\ \ \alpha V^B_{\ell j}S_{i\ell j}^{\T}\,]$
  needs their orthonormal bases, i.e.\ two full Householder QR factorizations
  of panels up to $b_m\times\hat r$ and $b_n\times\hat r$, on top of the
  $\hat r\times\hat r$ SVD used to rotate into the truncated basis.
\item \textbf{Memory traffic scales with $q_k$, not with the tile once.}
  Because the tile is never left implicit, its $O((b_m+b_n)r)$ factor pair is
  read out of and written back to device memory at every one of the $q_k$
  steps rather than once, and because rank grows by up to $r_{A,i\ell}$ before
  each truncation, the QR/SVD themselves run on panels wider than the
  converged rank $r_{C,ij}$ would require.
\end{itemize}

This is the metric the rest of the paper is built against: not the arithmetic
operation count alone, but the number and width of compressions and the
amount of factor-pair traffic they force.

\subsection{Randomized sketching: trading QR+SVD for GEMM}
\label{sec:randomized}

The randomized alternative developed here keeps every $A_{i\ell}B_{\ell j}$
term of \eqref{eq:update} \emph{implicit} for the whole reduction over $\ell$
and applies a blocked adaptive randomized range finder (ARA) to the entire sum
at once, so that a tile is compressed exactly once, at the very end,
regardless of $q_k$. Two changes follow from that decision:

\begin{enumerate}
\item \textbf{The QR/SVD-based recompression at every $\ell$ disappears.}
  Nothing about the tile is ever materialized as an explicit factor pair
  before convergence, so there is nothing to recompress at intermediate
  steps. The only factorizations are the range finder's Cholesky-QR
  normalization of a narrow sketch block (\S\ref{sec:cholqr}) --- a
  GEMM-plus-small-triangular-solve operation, not a Householder QR --- and,
  once, the final thin SVD of the co-range at \S\ref{sec:ara}'s exit. Total
  compressions per tile: one, not $q_k$.
\item \textbf{The reduction over $\ell$ becomes an ordinary GEMM contraction
  axis.} Rather than accumulating and re-truncating $q_k$ times, the sum
  $\sum_\ell A_{i\ell}B_{\ell j}$ is sketched as a whole: sketch first and sum
  the (narrow) samples afterwards, instead of summing (wide) tiles and
  sketching the sum. \S\ref{sec:factorlist} makes this precise; the key fact
  is that a random sketch commutes with the outer reduction because matrix
  multiplication is bilinear, so the $\ell$ sum can be folded into the
  contraction dimension of one large batched GEMM instead of being unrolled
  into $q_k$ sequential steps.
\end{enumerate}

Both properties reduce, in the language of \S\ref{sec:setting}, to keeping the
working set proportional to $\hat\rho_{ij}$ and the sketch width $s$ rather
than to $b_m$, $b_n$, or $q_k$ directly --- quantified in \S\ref{sec:minws}
and \S\ref{sec:cost}.

\section{The tile update as a factor list}
\label{sec:factorlist}

Using the coupling matrix \eqref{eq:coupling}, the update \eqref{eq:update} is
a sum of $n_f$ outer products,
\begin{equation}
\label{eq:list}
X_{ij}=\sum_{a=1}^{n_f}L_aR_a^{\T},
\qquad
L_a\in\R^{b_m\times\nu_a},\quad R_a\in\R^{b_n\times\nu_a},
\end{equation}
where the coupling is absorbed into the narrower side so that
$\nu_\ell=\min(r_{A,i\ell},r_{B,\ell j})$,
\begin{equation}
\label{eq:absorb}
(L_\ell,R_\ell)=
\begin{cases}
\bigl(U^A_{i\ell},\;\alpha V^B_{\ell j}S_{i\ell j}^{\T}\bigr),
& r_{A,i\ell}\le r_{B,\ell j},\\[4pt]
\bigl(U^A_{i\ell}S_{i\ell j},\;\alpha V^B_{\ell j}\bigr),
& r_{A,i\ell}>r_{B,\ell j}.
\end{cases}
\end{equation}
The term $(L_C,R_C)=(U^C_{ij},\beta V^C_{ij})$ is included when $\beta\ne0$
(\S\ref{sec:beta}). Unlike \textsc{DeterministicUpdate}
(Algorithm~\ref{alg:kblas}), this list is never assembled, concatenated, or
recompressed as a running accumulator; it is an iterator whose pairs are
produced on demand, consumed, and discarded. Everything downstream is built
from two implicit operators,
\begin{equation}
\label{eq:apply}
\mathrm{ApplyRight}(\Omega)=X_{ij}\Omega=\sum_aL_a\bigl(R_a^{\T}\Omega\bigr),
\qquad
\mathrm{ApplyLeft}(Q)=X_{ij}^{\T}Q=\sum_aR_a\bigl(L_a^{\T}Q\bigr).
\end{equation}

\subsection{Sketching the right-hand side}
\label{sec:ordering}

Fix a sketch block $\Omega\in\R^{b_n\times s}$ with $s$ small (tens of
columns, versus $b_n$ in the hundreds). \eqref{eq:apply} already sums small
products, but it is worth deriving $\mathrm{ApplyRight}$ directly from
\eqref{eq:update} to see exactly what commutes with what. By distributivity
and associativity of matrix multiplication,
\begin{align}
X_{ij}\Omega
&=\Bigl(\alpha\sum_{\ell=1}^{q_k}A_{i\ell}B_{\ell j}+\beta C_{ij}\Bigr)\Omega
\notag\\
&=\alpha\sum_{\ell=1}^{q_k}U^A_{i\ell}S_{i\ell j}\bigl((V^B_{\ell j})^{\T}
\Omega\bigr)+\beta U^C_{ij}\bigl((V^C_{ij})^{\T}\Omega\bigr),
\label{eq:distrib}
\end{align}
substituting $A_{i\ell}B_{\ell j}=U^A_{i\ell}S_{i\ell j}(V^B_{\ell j})^{\T}$
term by term: the sketch is pushed onto the \emph{right-hand factor} of every
term before anything is summed. Two definitions turn \eqref{eq:distrib} into
a fused reduction. With
\begin{equation}
\label{eq:HT}
H_{\ell j}:=\bigl(V^B_{\ell j}\bigr)^{\T}\Omega\in\R^{r_{B,\ell j}\times s},
\qquad
T_{i\ell j}:=\alpha\,S_{i\ell j}H_{\ell j}\in\R^{r_{A,i\ell}\times s},
\end{equation}
the sample is
\begin{equation}
\label{eq:Y}
Y_{ij}=X_{ij}\Omega=\sum_{\ell=1}^{q_k}U^A_{i\ell}\,T_{i\ell j}
\;+\;\beta\,U^C_{ij}\bigl((V^C_{ij})^{\T}\Omega\bigr).
\end{equation}
Every quantity computed before the final sum in \eqref{eq:Y} is small:
$H_{\ell j}$ and $T_{i\ell j}$ each have one dimension equal to the sketch
width $s$ and the other equal to a rank, never to $b_m$ or $b_n$. For a
representative off-diagonal tile with rank $r\sim30$ and sketch width
$s\sim16$ against a tile edge $b\sim300$, $H_{\ell j}$ and $T_{i\ell j}$ hold
on the order of $rs\approx500$ entries each, against $b^2\approx9\times10^4$
for the dense tile that is never formed. The long dimension $b_m$ is touched
exactly once, by the final GEMM \eqref{eq:Y}.

That final GEMM is where the reduction over $\ell$ is folded into an ordinary
contraction axis rather than unrolled into $q_k$ sequential steps.
Concatenating the per-$\ell$ pieces horizontally and vertically,
\begin{equation}
\label{eq:foldedgemm}
\mathbf U_i:=\bigl[\,U^A_{i1}\ \cdots\ U^A_{iq_k}\,\bigr]\in\R^{b_m\times\rho_i},
\qquad
\mathbf T_{ij}:=\begin{bmatrix}T_{i1j}\\\vdots\\T_{iq_kj}\end{bmatrix}
\in\R^{\rho_i\times s},
\end{equation}
\eqref{eq:Y}'s sum (the $\beta$ term aside) is literally the single product
$Y_{ij}=\mathbf U_i\mathbf T_{ij}$: one GEMM with contraction dimension
$\rho_i$. \textsc{DeterministicUpdate} also concatenates factors (its
line~3) --- but it concatenates \emph{into a recompression}, at every step,
whereas \eqref{eq:foldedgemm} concatenates \emph{into a GEMM's contraction
axis}, once per sketch block, with no factorization in between. The
implementation never materializes $\mathbf U_i$ or $\mathbf T_{ij}$ as literal
contiguous arrays either: the same result comes from one batched-GEMM call
whose stride/batch structure already treats $\ell$ as an extra contraction
index, but \eqref{eq:foldedgemm} is the clearest way to see why this is one
GEMM and not $q_k$.

The alternative ordering $(\sum_aL_aR_a^{\T})\Omega$ would instead form a
$b_m\times b_n$ object. Explicitly, the two are algebraically identical,
\[
Y=\sum_a L_a\bigl(R_a^{\T}\Omega\bigr)
\quad\text{(inner size }\nu_a\times s)\qquad\text{versus}\qquad
Y=\sum_a\bigl(L_aR_a^{\T}\bigr)\Omega
\quad\text{(inner size }b_m\times b_n),
\]
but only the first has cost proportional to $\hat\rho_{ij}s$ rather than
$b_mb_n$. The transpose apply $X^{\T}Q$ is evaluated symmetrically,
eliminating $b_m$ first through $(U^A_{i\ell})^{\T}Q$.

\begin{remark}[What gets hoisted and reused]
\label{rem:hoist}
$S_{i\ell j}$ is independent of $\Omega$ and is computed once per tile,
retained across all range-finder iterations --- the analogue of
\textsc{DeterministicUpdate}'s $S_{i\ell j}$ (Algorithm~\ref{alg:kblas},
line~3), except there it is recomputed at whatever moment $\ell$ is visited
and consumed immediately into a recompression, while here it is computed once
and read many times. Only $H_{\ell j}$ and $T_{i\ell j}$ are regenerated per
sketch block. Moreover $H_{\ell j}$ carries no row index $i$: fixing a block
column $j$ and sweeping $i$ computes $H_{\ell j}$ once and reuses it across
the whole column. Nothing in \textsc{DeterministicUpdate} is reusable this
way --- its running accumulator $(U,V)$ is private to one output tile by
construction, since it is exactly the thing being recompressed. Column-wide
sharing of $H_{\ell j}$ is one of the two reasons the schedule of
\S\ref{sec:schedule} traverses by column (or, symmetrically, by row).
\end{remark}

\begin{remark}[Batch shape]
\label{rem:batch}
The association order in \eqref{eq:HT} is chosen per term from the ranks: the
coupling is contracted against whichever operand is cheaper. Because ranks vary
with $\ell$, the batch over $\ell$ in stages \eqref{eq:HT} is a grouped
(nonuniform) GEMM. The fused reduction \eqref{eq:Y}, by contrast, is a single
GEMM per tile with contraction dimension $\rho_i$: it is the large,
high-arithmetic-intensity kernel that dominates the useful work and that the
whole factored form exists to expose.
\end{remark}

\subsection{ApplyRight versus ApplyLeft}
\label{sec:applyroles}

The two operators of \eqref{eq:apply} play different roles inside the range
finder (Algorithm~\ref{alg:ara}) and are not interchangeable simplifications
of each other.

$\mathrm{ApplyRight}(\Omega_\Delta)=X_{ij}\Omega_\Delta$ is applied to a
\emph{fresh Gaussian} sketch every pass. Its output $Y_\Delta$ is what gets
orthogonalized against the existing basis and tested by the stopping rule
(\S\ref{sec:samplestop}): it is the mechanism that discovers new range
directions, and it is why the norm estimate \eqref{eq:normest} is unbiased
--- the randomness lives entirely in $\Omega_\Delta$.

$\mathrm{ApplyLeft}(Q_\Delta)=X_{ij}^{\T}Q_\Delta$ is applied to the
\emph{deterministic, already-orthonormalized} increment $Q_\Delta$ produced
by that same pass, not to a random sketch. Its output $Z_\Delta$ accumulates
the co-range coefficients the final truncation (\S\ref{sec:ara}) needs:
$Z=X_{ij}^{\T}Q$ is exactly the matrix whose thin SVD gives the optimal
$U,V$ once $Q$ is judged converged. ApplyLeft never decides whether to stop;
it only records, for the directions ApplyRight already confirmed, their
coordinates against the operator.

Both share the same cost structure and the same hoisting: ApplyLeft's
reduction eliminates $b_m$ first through $(U^A_{i\ell})^{\T}Q_\Delta$ exactly
as \eqref{eq:HT} eliminates it for ApplyRight, and its analogue of
$H_{\ell j}$ is equally column- (or row-) shared. The asymmetry is only in
what feeds them: random noise on the right, confirmed basis vectors on the
left.

\subsection{Canonical storage and sampling contract}
\label{sec:canonicalcontract}

The first implementation deliberately exposes a predictable restricted API:
$C$, $A$, and $B$ are all stored physically tile-row-major, while the transpose
flags define only their logical orientations.  On a regular tile grid this
gives the following operation-level schedule:
\begin{center}
\begin{tabular}{llll}
\toprule
$\operatorname{op}(A)$ & $\operatorname{op}(B)$ & repeated sample & run family \\
\midrule
$N$ & $N$ & $X\Omega$ & fixed output row \\
$N$ & $T$ & $X\Omega$ or $X^{\T}\Omega$ & fixed column or fixed row \\
$T$ & $T$ & $X^{\T}\Omega$ & fixed output row \\
$T$ & $N$ & deferred & packed/reduced general-storage path \\
\bottomrule
\end{tabular}
\end{center}
In the $NT$ case both terminal contraction stacks are natural.  Their retained
workspace is estimated from the active rank caps of the two operands, and one
side is selected for the complete operation.  Selection is not made per tile:
one run family preserves uniform batches and permits converged members to be
swap-removed into a dense active prefix.

``Natural'' here concerns both tile order and rank extent.  A panel at its
physical rank capacity is a zero-copy view.  If metadata trims that capacity,
the active rank prefix has gaps between adjacent tiles and cannot be reshaped
into the terminal GPU GEMM.  It is therefore packed once into compact
run-local storage and reused by all range-finder passes.  The more general
$TN$ packing/reduction problem is intentionally postponed rather than hidden
inside this canonical contract.

\section{Stopping: what is computed, and what it costs}
\label{sec:energy}

The loop stops on a quantity it has already computed, and it does not verify the
result afterwards. This section gives the test, then the two reasons --- one
statistical, one economic --- that no exact residual is formed on the hot path.

\subsection{The sampling stop}
\label{sec:samplestop}

After a block is projected and orthonormalized, the norms of the freshly
projected columns are already available: they are the diagonal of the Cholesky
factor $R$ produced by the intra-block orthonormalization. The classical
range-finder rule stops when a prescribed number $p$ of consecutive projected
columns have norm below the tolerance scale, i.e.\ when successive blocks stop
adding energy to the basis. Concretely, with
$a_e=\max_{\text{last }p}\|Y_\Delta(:,\cdot)\|$ and $R_{\max}$ the running
maximum of these norms, the loop stops when $a_e/R_{\max}$ falls below the
relative tolerance. This costs an $\mathcal O(s)$ reduction over numbers the
Cholesky QR already produced: no extra GEMM and no extra pass.

The diagonal is the right signal, and column norms are not. At convergence the
block is a set of random combinations of the few directions that remain, so
every \emph{column} still has $\mathcal O(1)$ norm while the \emph{block} is
rank deficient. By contrast $R[j,j]$ is the residual of column $j$ against both
the existing basis and the earlier columns of the same block, so it collapses
exactly when there is nothing new left.

An unbiased norm estimate is also free from the first sample: for a Gaussian
$\Omega\in\R^{b_n\times s}$, $\mathbb E\|X\Omega\|_F^2=s\|X\|_F^2$, so
\begin{equation}
\label{eq:normest}
\widehat{\|X\|}_F=\|Y\|_F/\sqrt{s}
\end{equation}
serves the relative tolerance $\varepsilon_{\mathrm{rel}}\|X\|_F$ at the cost of
one $\mathcal O(b_ms)$ norm.

\subsection{Why no exact verification}
\label{sec:noverify}

\paragraph{The stop is a bound, not a heuristic.}
Because each block is drawn from a \emph{fresh} Gaussian, $a_e$ is the
Halko--Martinsson--Tropp a posteriori estimator of
$\|(I-QQ^{\T})X\|_2$, whose failure probability decays like $10^{-p}$ in the
number of consecutive small blocks. The confidence is therefore a tuning knob:
raising $p$ by one costs a single block and buys an order of magnitude. Any
scheme that instead \emph{verifies} must beat that exchange rate.

\paragraph{It cannot.}
Forming $\Fnorm{X}^2$ from the factor list costs $\mathcal O(n_f^2)$ small
Grams, i.e.\ $\mathcal O(b_m\,q_k^2r^2)$ for representative rank $r$, whereas one
further sketch block costs $\mathcal O(b_m\,q_k r\,\Delta s)$. The ratio is
$q_k r/\Delta s$, about $32$ at $q_k=16$, $r=32$, $\Delta s=16$: it is far
cheaper to draw another block than to check whether one was needed. Verification
would only pay if it usually answered ``stop'', but the only situation that
would invoke it is one in which the stop already looked doubtful.

\paragraph{And the rank cap needs no residual either.}
If $r_{\max}$ is reached before convergence, that fact is recorded in the output
rank vector at no cost --- a tile returned at full rank is self-evidently
capped. Moreover the \emph{truncation} part of the error is exactly
$\sum_{k>r}\lambda_k$, available from the eigenvalues already computed for the
final rotation (\S\ref{sec:ara}); only the range-capture part is unknown, and
under a cap the truncation term dominates. So the honest error report is free.

\begin{remark}[Verification as a test-only facility]
\label{rem:verify}
The exact Frobenius residual is still computable from the factor list without a
dense tile, via $\Fnorm{X}^2=\sum_{a,b}\langle L_a^{\T}L_b,R_a^{\T}R_b\rangle_F$
and the analogous cross term. It is worth implementing behind a debug flag, as
an oracle for tests of the sampling stop. It should not appear on the execution
path, and it carries a cancellation floor of its own: it subtracts quantities of
size $\Fnorm{X}^2$ to obtain a residual of size $\tau^2$, so it is meaningful
only for $\tau^2$ safely above $u_{\mathrm{hi}}\Fnorm{X}^2$.
\end{remark}

\section{Folding in the existing \texorpdfstring{$\beta C_{ij}$}{beta C} tile}
\label{sec:beta}

The output tile $C_{ij}$ already exists in the format, as a low-rank pair
$(U^C_{ij},V^C_{ij})$ (or, on the diagonal, as a dense block). The GEMM
overwrites it with the accumulated $X_{ij}$ of \eqref{eq:update}. There are two
ways to include the $\beta C_{ij}$ term; we adopt the first and explain why.

\paragraph{As one further factor pair (adopted).}
Append $(L_C,R_C)=(U^C_{ij},\,\beta V^C_{ij})$ to the factor list. Then the range
finder builds a basis for the \emph{entire} update $X_{ij}$, product terms and
old tile together, and a single compression is performed. Concretely the term
enters the sample \eqref{eq:Y} as the extra summand
$\beta U^C_{ij}\bigl((V^C_{ij})^{\T}\Omega\bigr)$, and the transpose apply gains
$\beta V^C_{ij}\bigl((U^C_{ij})^{\T}Q\bigr)$. Its width $r_{C,ij}$ adds to $\hat\rho_{ij}$. When $C_{ij}$ is a dense diagonal block $D$, the pair is
replaced by the operators
$\Omega\mapsto\beta D\Omega$ and $Q\mapsto\beta D^{\T}Q$ directly; no
factorization of $D$ is needed for the apply.

\paragraph{As a post-merge (rejected).}
One could instead compress $\alpha\sum_\ell A_{i\ell}B_{\ell j}$ alone to a pair
$(Q,Z)$ and afterwards merge $[\,U^C_{ij}\ \ Q\,]$, recompressing an object of
width $r_{C,ij}+s$. This performs two truncations rather than one and introduces
a second, avoidable approximation of the old tile --- the same failure mode, at
smaller scale, as \textsc{DeterministicUpdate}'s per-step recompression
(\S\ref{sec:kblas}). Carrying $\beta C_{ij}$ inside the range finder makes $Q$
a range basis for the complete update and truncates exactly once.

\begin{remark}[Why the format's orthonormality is not enough for a norm]
\label{rem:betanorm}
Because the update mixes the product terms with the stored tile, $X_{ij}$ has no
closed-form Frobenius energy from the format's orthonormal-basis invariant alone.
This costs nothing here: the sampling estimate \eqref{eq:normest} requires no
closed form, and it is the only norm the execution path uses.
\end{remark}

\section{The blocked randomized range finder}
\label{sec:ara}

This is the range finder that replaces \textsc{DeterministicUpdate}
(Algorithm~\ref{alg:kblas}): it produces the same low-rank output but performs
exactly one compression per tile, never one per contraction step
(\S\ref{sec:randomized}). A trial basis $Q\in\R^{b_m\times s_Q}$ for the
column space of $X_{ij}$ is grown one block at a time; the projection
$Z=X^{\T}Q$ is accumulated in step. The loop stops on the free sampling
estimate of \S\ref{sec:samplestop} and does not verify afterwards
(\S\ref{sec:noverify}). Block orthogonalization is classical block
Gram--Schmidt with one reorthogonalization step (BCGS2); the intra-block
orthonormalization is Cholesky QR, with the Gram and its Cholesky factor computed
in a promoted precision while data is stored in the working precision.

\begin{algorithm}[H]
\caption{\textsc{RangeFind} --- compress one factor list to tolerance}
\label{alg:ara}
\begin{algorithmic}[1]
\Procedure{RangeFind}{$\{(L_a,R_a)\},\,\varepsilon_{\mathrm{rel}},\varepsilon_{\mathrm{abs}},r_{\max},s,\Delta s,s_{\max}$}
  \State $Q\gets[\,]$, $Z\gets[\,]$, $s_Q\gets0$
  \While{$s_Q<s_{\max}$}
     \State draw Gaussian $\Omega_\Delta\in\R^{b_n\times d}$,
            $d=\min(\Delta s,\,s_{\max}-s_Q)$ (or $s$ on the first pass)
     \State $Y_\Delta\gets\mathrm{ApplyRight}(\Omega_\Delta)$
        \Comment{fused sample \eqref{eq:Y}}
     \If{$s_Q=0$}\ $\widehat{\|X\|}_F\gets\|Y_\Delta\|_F/\sqrt d$;\
        $\tau^2\gets\max(\varepsilon_{\mathrm{abs}}^2,
        \varepsilon_{\mathrm{rel}}^2\widehat{\|X\|}_F^2)$
     \EndIf
        \Comment{norm estimate \eqref{eq:normest}, free}
     \State $Q_\Delta\gets\Call{OrthAgainst}{Q,Y_\Delta}$
        \Comment{BCGS2 $+$ Cholesky QR; column norms $=\mathrm{diag}(R)$}
     \State update $a_e,R_{\max}$ from the projected column norms
        \Comment{sampling stop, \S\ref{sec:samplestop}}
     \State $Z_\Delta\gets\mathrm{ApplyLeft}(Q_\Delta)$
        \Comment{fused $X^{\T}Q$, \S\ref{sec:applyroles}}
     \State $Q\gets[\,Q\ \ Q_\Delta\,]$;\ $Z\gets[\,Z\ \ Z_\Delta\,]$;\
            $s_Q\gets\mathrm{cols}(Q)$
     \If{$p$ consecutive projected columns satisfy $a_e/R_{\max}<\varepsilon_{\mathrm{rel}}$}
        \State \textbf{break}
           \Comment{\S\ref{sec:noverify}: no verification pass}
     \EndIf
  \EndWhile
  \State $K\gets Z^{\T}Z$ (promoted);\
         $(\lambda,W)\gets\textsc{XsyevBatched}(K)$; reverse to descending
  \State $r\gets\min\bigl(r_{\max},\ \min\{t:\textstyle\sum_{k>t}\lambda_k\le\tau^2\}\bigr)$
     \Comment{$E_{\mathrm{pred}}$ screen}
  \State $U\gets QW(:,1{:}r)$;\ $V\gets ZW(:,1{:}r)$
  \State \Return $(U,\,V,\,\textstyle\sum_{k>r}\lambda_k,\;[\,r=r_{\max}\,])$
     \Comment{achieved truncation error and the cap flag, both free}
\EndProcedure
\end{algorithmic}
\end{algorithm}

There is a single exit. The sampling estimate decides when to stop and the tail
sum $\sum_{k>r}\lambda_k$ against $\tau^2$ chooses the rank; both are byproducts
of work already done, so the routine has no branch that costs an extra pass over
the factor list. Randomness affects the pass count and, at the probability of
\S\ref{sec:noverify}, the achieved accuracy.

The final eigendecomposition is not a numerical crutch --- $Z$ is not rank
deficient by construction --- but it does real work: blocked range finders
overshoot the true rank by a few columns, since sampling proceeds in units of
$\Delta s$, and rotating into the eigenbasis of $K$ before truncating recovers
the optimal rank for the basis that was built.

\subsection{Why Cholesky QR is safe here}
\label{sec:cholqr}

Cholesky QR squares the condition number and is unstable on rank-deficient input.
It is nonetheless the right orthonormalizer inside this loop, for a structural
reason. The block handed to it has already been projected against the existing
basis by BCGS2, so in the normal case it is well conditioned and mixed-precision
Cholesky QR (Gram and factor in promoted precision, data in working precision)
suffices. The dangerous case is the final block at convergence, whose columns
nearly lie in $\mathrm{span}(Q)$; those columns then have small norm and are
flagged as converged, i.e.\ discarded, rather than trusted as basis vectors. The
instability therefore lands precisely on the vectors that are thrown away. This
safety depends on keeping $s$ small: enlarging the block to cut the pass count
re-exposes it to ill conditioning. The large-GEMM efficiency should be recovered
from the fused reduction \eqref{eq:Y}, whose contraction dimension is $\rho_i$
independent of $s$, not from widening $s$.

\subsection{Primitive inventory}
\label{sec:primitives}

Every kernel in Algorithm~\ref{alg:ara} is a stock batched routine:
the applies \eqref{eq:HT}--\eqref{eq:Y} and the projections are batched
(grouped) GEMM; the Gram is batched SYRK in promoted precision; the Cholesky is
batched POTRF; the back-solve is batched TRSM; the small $s_Q\times s_Q$
eigenproblem is \texttt{XsyevBatched} (exact, uncapped in the wrapper). No
routine that forms an explicit Householder $Q$ factor is required, so no
dependence on a library outside cuBLAS/cuSOLVER is incurred --- a direct
contrast with Algorithm~\ref{alg:kblas}'s two Householder QRs per contraction
step (\S\ref{sec:kblas}). The only custom kernels are
small reductions: the projected-column norms feeding the stopping test, and the
per-member convergence bookkeeping that decides how many columns each tile of a
batch samples next.

\section{Scheduling and workspace}
\label{sec:schedule}

The implemented scheduler preserves the sharing that determines the canonical
sampling family and then rolls work through a reusable lane. A lane is either
a fixed output column or a fixed output row. A fixed-column
$\texttt{ColumnRunCoupling}$ forms the shared $H_{\ell j}$ once, while the fixed-row
couplings retain their corresponding shared operand. Arbitrary output tiles
are not admitted into one larger lane: doing so would discard this sharing and
require a new mixed-tile coupling.

The first implementation is deliberately single-lane, lane-local, and
full-$k$. The workspace determines the number of simultaneously active members
in that lane, not a rectangular macrotile or a new contiguous run per pass.
Members that finish are retired in waves and their slots are reused by pending
output tiles.

\subsection{Minimum workspace per unit of work}
\label{sec:minws}

The working set for one tile in flight, at sketch width $s$, is
\begin{equation}
\label{eq:tilews}
w_{ij}(s)=
\underbrace{(b_m+b_n)\,s_{\max}}_{\text{accumulators }Q,\,Z}
\;+\;
\underbrace{\textstyle\sum_\ell r_{A,i\ell}r_{B,\ell j}}_{\text{retained }S_{i\ell j}}
\;+\;
\underbrace{\rho_i\,s}_{\text{applied }T\text{ column}},
\end{equation}
sized at the cap $s_{\max}$ for the accumulators, since $Q$ and $Z$ grow to at
most that width. No $b_m\times b_n$ buffer ever appears --- nor does
\textsc{DeterministicUpdate}'s running $(U,V)$ pair need to be widened past its
converged rank at any point, since there is no recompression step to size for.
For a block column $j$, the coupling matrices $H_{\ell j}$ are shared, adding a
one-time $\sum_\ell r_{B,\ell j}\,s=\gamma_j s$ that is amortized over the
column rather than paid per tile. The minimum viable workspace to process a
single tile of the column is thus $w_{ij}(s_0)$ at the initial width $s_0$; the
minimum to run the column with maximal accumulator width is
$\max_i w_{ij}(s_{\max})+\gamma_j s_{\max}$.

These quantities are computable \emph{before any arithmetic} from rank metadata:
$\rho_i$, $\gamma_j$, and $\sum_\ell r_{A,i\ell}r_{B,\ell j}$ are all read from
the stored ranks. They describe admission and sampling cost, while the
implementation chooses capacity through exact typed arena-size formulas.

The arena separates whole-run persistent storage from phase storage. Persistent
storage contains the basis, packed factor stacks, and run state. Phase storage
is reused between admission, sampling, and retirement:
\begin{equation}
\label{eq:arena}
w_{\mathrm{run}}=w_{\mathrm{persistent}}+
\max\{w_{\mathrm{admission}},w_{\mathrm{sampling}},w_{\mathrm{retirement}}\}.
\end{equation}
Admission-time $S$ formation is released before sampling, and the same phase
allocation is reused for co-range application, truncation, and scatter. The
minimum workspace is the one-slot bound; the maximum is the complete-lane
bound. An intermediate byte budget selects the largest fitting slot capacity,
while a budget below the minimum is rejected and excess bytes are capped.

\subsection{Rolling traversal within one lane}
\label{sec:traversal}

For a lane with capacity $c$, the scheduler maintains an active prefix of at
most $c$ members and a pending queue. It admits an initial wave, forms each
member's one-time $S$ core over the complete contraction range, and clears the
slot-local ARA state. Every subsequent sampling pass operates on the active
prefix. Members converging at the same full-$k$ boundary form one retirement
wave; the scheduler applies the co-range, truncates, and scatters that wave
before admitting replacements into the released slots. A replacement joins
the current block schedule with fresh basis and convergence state; it does not
replay earlier sketch widths.

\begin{algorithm}[H]
\caption{\textsc{Schedule} --- rolling admission in one lane}
\label{alg:sched}
\begin{algorithmic}[1]
\Procedure{Schedule}{lane,\,$c$}
  \State pending $\gets$ output members of lane
  \State active $\gets$ \Call{AdmitWave}{first $\min(c,|$pending$|)$ members}
  \While{active $\neq\varnothing$}
     \State sample all active members through one complete full-$k$ pass
     \State $W\gets$ members converged at this boundary
     \If{$W\neq\varnothing$}
        \State \Call{RetireWave}{co-range, truncate, scatter, $W$}
        \State remove $W$ from active by swap-compaction
        \State admit up to $|W|$ pending members into released slots
        \State reset their basis, diagnostics, and convergence state
     \EndIf
  \EndWhile
\EndProcedure
\end{algorithmic}
\end{algorithm}

The active prefix is kept packed. Full-width members remain in front of a
terminal-width suffix, so a pass uses no more than two orthogonalization groups
while retaining one fused full-$k$ contraction. Retirement is wave-batched
rather than member-at-a-time; consequently co-range and finalization launch
growth is bounded by convergence boundaries, not by the number of output tiles.
Sampling parameters are cohort-global. This can oversample an easy late
arrival, so schedule statistics record sampled columns and FLOPs per member in
addition to active and pending widths.

\subsection{Reduction staging is deferred}
\label{sec:staging}

The current scheduler always retains the full contraction range in one fused
sampling operation. If a single tile's full-$k$ working set does not fit, the
operation fails its workspace contract rather than silently changing the
reduction. Chunking $k$ into panels would accumulate
$Y^{(K_t)}=\sum_{\ell\in K_t}U^A_{i\ell}T_{i\ell j}$ repeatedly and move roughly
$(2\kappa-1)b_ms$ sample elements for $\kappa$ panels. It remains a future
optimization, to be enabled only if full-$k$ transient storage, rather than
lane tails, limits occupancy.

\subsection{Workspace and deferred scheduling}
\label{sec:choose}

The public operation receives an optional reusable workspace. Internally,
$\texttt{CompressedGemmWorkspace}$ owns the arena, slot-local progress, convergence
and diagnostic arrays, output panels, and global scatter buffers. The workspace
compatibility key includes the selected sampling family, geometry, active rank
caps, block width, output rank cap, and slot capacity. The driver can therefore
reuse the same allocation across complete GEMM calls without changing the
schedule. Internal schedule statistics expose admissions, retirement waves,
active and pending widths, passes, padded projection columns, discarded
terminal columns, and phase timings without adding profiling keywords to the
public GEMM API.

The sketch width $s$ (and its cap $s_{\max}$) remains small enough that
Cholesky QR stays well conditioned (\S\ref{sec:cholqr}) and the accumulators
remain modest. The large-GEMM efficiency still comes from the fused reduction
\eqref{eq:Y}, whose contraction dimension $\rho_i$ is set by the data.

The next occupancy step is R4.5: combine ready members from several lanes into
one larger pointer-batched GEMM call. This recovers cross-lane occupancy without
independent streams or shared free-list state, but gives up some same-axis
sharing. A shared multi-lane arena and an arbitrary mixed-tile cohort are
separate deferred designs; the latter requires a new coupling. The $TN$
operation, arbitrary physical tile orders, boundary tiles, and reduction
sub-panelling likewise remain outside the canonical scheduling contract.

\section{Cost}
\label{sec:cost}

Per tile, to leading order and excluding the coupling prologue $S_{i\ell j}$
(common to any method):

\begin{center}
\begin{tabular}{lll}
\toprule
Stage & Cost & Note \\
\midrule
$H_{\ell j}$ (per sample)        & $2b_n\gamma_j s$ & once per column, $/q_m$ per tile \\
$Y_{ij}=X\Omega$                 & $2b_m\rho_i s$   & one fused GEMM, contraction $\rho_i$ \\
$Z_{ij}=X^{\T}Q$                 & $2b_m\rho_i s+2b_n\gamma_j s$ & two fused GEMMs \\
Gram $+$ eig (\textsc{XsyevBatched}) & $2b_ms_Q^2+O(s_Q^3)$ & promoted Gram, small eig \\
form $U,V$                       & $2(b_m+b_n)s_QR$ & final lift to output rank $R$ \\
norm estimate \eqref{eq:normest} & $O(b_ms)$ & free, every tile \\
\bottomrule
\end{tabular}
\end{center}

The number of small dense factorizations is $O(q_mq_n)$ (one eigenproblem per
tile, of size $s_Q$) against Algorithm~\ref{alg:kblas}'s $O(q_kq_mq_n)$, and
the chain of dependent launches per tile is $O(1)$:
the reduction over $\ell$ is a single fused GEMM, not a serial recurrence. Batch
shapes at fixed $s$ are uniform across a run, and the number of range-finder
passes is the only data-dependent quantity. No term is quadratic in $q_k$: the reduction is never
revisited once sampled, which is what removing the verification pass buys.

\section{Summary of the blueprint}
\label{sec:summary}

The baseline is a per-tile blocked range finder (Algorithm~\ref{alg:ara}) driven
by the rolling single-lane scheduler (Algorithm~\ref{alg:sched}), replacing the
deterministic sequential recompression of Algorithm~\ref{alg:kblas}
(\S\ref{sec:strategies}). The
reduction over $\ell$ is folded into a single high-intensity GEMM
\eqref{eq:Y} by the long-dimension-last ordering; the existing $\beta C_{ij}$
tile is carried inside the range finder as one further factor pair
(\S\ref{sec:beta}) so that a single truncation produces the output; the loop
stops on the free sampling estimate (\S\ref{sec:samplestop}), which is an a
posteriori bound rather than a heuristic, and never verifies afterwards, because
drawing another block is roughly $q_kr/\Delta s$ times cheaper than checking
whether one was needed (\S\ref{sec:noverify}); and every kernel is a
stock batched cuBLAS/cuSOLVER routine (\S\ref{sec:primitives}), with Cholesky QR
made safe by the convergence-time discarding of small columns
(\S\ref{sec:cholqr}). Rank metadata determines the lane family and workspace
capacity before arithmetic; rolling admission and retirement then handle the
data-dependent convergence without reallocating the numerical arena. The net
effect against \S\ref{sec:kblas}'s baseline is a compression count cut from
$O(q_kq_mq_n)$ to $O(q_mq_n)$ and every recompression's Householder QR
replaced by batched GEMM and Cholesky-QR. Full-$k$
reduction remains the baseline; multi-lane growth and reduction staging are
deferred.

\end{document}
